\documentclass[zad,zawodnik,utf8,pl]{sinol}
\title{TODO}
\id{__ID__}
\signature{TODO}
\konkurs{TODO}
\day{TODO}
\date{TODO}
\RAM{512}
\iomode{stdin}
\pagestyle{fancy}

\usepackage{amsmath}
% TEX_PACKAGE_GUARD
% Remember to add the package to the github action whenever using
% packages not present in TeXLive small.

% TEMPLATE: the layout of the library tasks from the camps. The tests are encrypted, so
% \makecompactexample would print garbage: show the example as a table of calls.

\begin{document}
\begin{tasktext}
	\vspace{-0.5cm}
	\noindent {}

	\noindent TODO: historyjka i zadanie. \\

	\section{Komunikacja}
	To zadanie jest z biblioteczką. Aby jej użyć, należy wpisać na początku programu:
	\begin{itemize}
        \item \textbf{C++:}
            \verb|#include "__ID__lib.h"|
    \end{itemize}
	Biblioteka udostępnia następujące funkcje:
	\begin{itemize}
        \item $\texttt{dlugosc}()$ \\
            Zwraca $n$ -- TODO.
            \begin{itemize}
                \item \textbf{C++:}
		                \texttt{int dlugosc();}
            \end{itemize}
        \item $\texttt{element}(i)$ \\
            Zwraca $a_i$ ($1 \leq i \leq n$) -- TODO.
            \begin{itemize}
                \item \textbf{C++:}
		                \texttt{long long element(int i);}
            \end{itemize}
    \end{itemize}
	Twój program musi zawierać następującą funkcję:
	\begin{itemize}
        \item $\texttt{maks\_suma}()$ \\
            Biblioteka wywoła ją dokładnie raz. Funkcja ma zwrócić TODO.
            \begin{itemize}
                \item \textbf{C++:}
		                \texttt{long long maks\_suma();}
            \end{itemize}
    \end{itemize}

	\noindent Twój program \textbf{nie może} czytać żadnych danych
    (ani ze standardowego wejścia, ani z plików).
    \textbf{Nie może} również nic wypisywać do plików ani na standardowe wyjście
    ani zawierać funkcji \texttt{main}.
    Może pisać na standardowe wyjście diagnostyczne
    (\texttt{stderr}) -- pamiętaj jednak, że zużywa to cenny czas.\\

	\section{Założenia}
	\begin{itemize}
		\item $1 \leq n \leq 200\,000$
		\item $|a_i| \leq 10^9$
	\end{itemize}

	\section{Ocenianie}
	\begin{itemize}
		\item Funkcję \texttt{element($i$)} można wywołać tylko dla $1 \leq i \leq n$ i co najwyżej dwa razy dla tego samego $i$. W przeciwnym razie program automatycznie się kończy i otrzymuje $0$ punktów za test.
		\item Za błędną odpowiedź otrzymasz $0$ punktów za test.
		\item W przypadku błędów z komunikacją twój program otrzyma werdykt \textit{Błąd związany z komunikacją}.
	\end{itemize}

	\section{Przykładowy przebieg programu}

	\begin{center}
		\begin{tabular}{|l|l|p{10cm}|}
			\hline
			\textbf{Wywołana funkcja} & \textbf{Wynik}  & \textbf{Opis} \\ \hline
			\texttt{maks\_suma()}   &         & Biblioteka wywołuje funkcję zawodnika. \\ \hline
			\texttt{dlugosc()}      & $3$     & TODO \\ \hline
			\texttt{element($1$)}   & $1$     & TODO \\ \hline
			                        & TODO    & Funkcja \texttt{maks\_suma} zwraca TODO. \\ \hline
		\end{tabular}
	\end{center}

	\section{Podzadania}
	% Keep in sync with `scores` and `time_limit(s)` in config.yml (checked by `contest.py lint`).
	\begin{center}
		\begin{tabular}{|l|l|p{7cm}|l|}
			\hline
			\textbf{Podzadanie} & \textbf{Liczba punktów} & \textbf{Ograniczenia} & \textbf{Limit czasu} \\ \hline
			1 & $30$ & $n \leq 1000$ & $1$ s \\ \hline
			2 & $70$ & brak dodatkowych założeń & $1$ s \\ \hline
		\end{tabular}
	\end{center}

	\section{Biblioteka}

	\noindent Możesz założyć, że działanie biblioteczki, jeżeli wszystkie funkcje zostaną wywołane maksymalną liczbę razy, nie zajmie więcej niż TODO s oraz nie zajmie więcej niż TODO MB pamięci RAM.

	\section{Eksperymenty}

	\noindent{\sloppy Przykładowe \textbf{błędne} rozwiązanie i przykładowa biblioteka w języku C++ znajdują się w archiwum \texttt{dlazaw.zip} w dziale \textit{Pliki i testy} na SIO. Rozwiązanie to i biblioteka obrazują typy zwracane oraz przyjmowane przez wszystkie funkcje. Przykładowa biblioteka może różnić się zachowaniem od tej, która zostanie użyta do oceny rozwiązań, i może nie spełniać założeń zadania. Ma ona jedynie pokazać sposób interakcji z programem. Wczytuje ze standardowego wejścia TODO. Przykładowe wejście znajduje się w pliku \texttt{__ID__0.in}.\par}

  \noindent Rozwiązanie \texttt{__ID__.cpp} w języku C++ można skompilować i uruchomić w następujący sposób:

  \texttt{g++ -O3 -static -o __ID__ __ID__.cpp __ID__lib.cpp -std=c++23}

  \texttt{./__ID__ < __ID__0.in}

\noindent Należy zadbać o to, aby pliki \texttt{__ID__lib.cpp} oraz \texttt{__ID__lib.h} znajdowały się w tym samym folderze co rozwiązanie.

\end{tasktext}
\end{document}
